\chapter{Introduction}
With more \gls{uav}s in the air each year, and a growing desire 
for them to operate in controlled airspaces, there is a concern regarding their ability 
to maintain separation from hazards in the air. United States and European federal 
agencies are developing policy that would require \gls{daa} systems for
\gls{uav}s, which parallel the \gls{vfr} and \gls{ifr} that dictate
how human pilots avoid traffic \cite{easa_easy_access_uas_2024} \cite{rtca_sc228_2025} 
\cite{serres_wu_2021_nasa_tm}.

The extensive research on \gls{eo} detection systems for UAVs, 
demonstrates the high angular resolution, and recognition capabilities 
of \gls{eo} sensors in well-lit conditions \cite{rs16050756}, but also highlights 
their susceptibility to low-light, 
precipitation, and motion blur \cite{drones8110638}. While there are 
techniques to mitigate these issues, such as through the use of infrared 
cameras or  
image enhancement techniques \cite{aerospace9120829}, they provide limited 
performance improvements in harsh real-world weather conditions. Another
important shortcoming of \gls{eo} sensors is that they do not naturally provide
range information, which is critical for collision avoidance.

Another category of sensors for UAV DAA is LiDAR, which creates precise 
3D maps based on returns from laser pulses. As an active sensor, LiDAR 
performs better than cameras in low light. However, surveys   
identify the susceptibility of LiDAR to strong light sources in the 
environment, and weather conditions such as snow, rain and fog \cite{ZHANG2023146}.
Furthermore, multibeam LiDAR systems designed for small UAVs often sacrifice 
performance, limiting lightweight models to ranges of 100--200 m 
or low frame rates \cite{benewake_tfa1200}  
\cite{sysak2025_lightweight_compact_pulse_radar_uav_collision_avoidance}. 
While improvements to range and resolution continue, the fundamental 
limitation of LiDAR motivates the exploration of a more all-weather capable
modality. 

Radar sensors may be the most suitable for \gls{daa}, as they can perform 
sensing over great distances in all weather conditions, due to their relatively low 
atmospheric attenuation \cite{borden_radar_imaging_airborne_targets_1999}. The 
inconvenient side effect of their long wavelength property is that angular resolution 
is much lower compared to similarly sized optical systems. If the challenges stemming 
from this limitation are overcome, the implementation of radar as the primary sensing 
mode or an additional sensing mode onto \gls{uav}s will improve their perception capabilities 
and safety.

Historically, airborne radar has been used for large aircraft, where size weight and 
power (SWaP) constraints are not strict. The size and weight of radar systems has been
the largest barrier to their use on smaller platforms, \cite{Lies2020}, \cite{McLain2015},
\cite{Wang2016}. This historical constraint is now weakening because both sides of the
radar stack have miniaturised. Texas Instrument's AWR2243 is a single 76-81 GHz radar chip
that integrates \gls{pll}s, three transmitters, four receivers, a mixer, an
\gls{adc}, and self calibration, all on a 10.4mm $\times$ 10.4mm \gls{pcb}-friendly package. Higher frequency, lower noise, electronics and the miniaturisation of 
the transistor have been key enablers of this technology. Higher frequency radar allows for
higher resolution imaging or smaller aperture; lower noise allows for increased
sensitivity to targets. On the compute side, modern UAV systems often have a dedicated, 
high compute companion computer in addition to the flight compute. NVIDIA’s Jetson series
computers such as Xavier NX, Orin Nano and Orin NX provide 21, 67, and 157 \gls{tops}
on 8-bit integers respectively, while each weigh under
50g (without a heat sink).

Plextek's mmWave-based 60 GHz sense-and-avoid radar for \gls{uav}s demonstrates the feasibility
of pairing miniaturised radar hardware with modern edge computing for the detection of
people, vehicles, poles, and buildings at ranges of hundreds of metres \cite{Plextek2024}. Echodyne’s
EchoFlight, a 817 g commercial radar systems for \gls{uav}s that can detect small drone targets
at a range of 750 m, reflects the onset of the technology into industry \cite{Echodyne2024}. The wider \gls{uav} market shows a similar shift toward more computationally
demanding onboard tasks. Skydio’s X10, for instance, incorporates both an NVIDIA Jetson
Orin \gls{soc} and a Qualcomm QRB5165 \gls{soc} within a 2.11 kg platform while retaining up to 40
minutes of flight time \cite{Skydio2023}, illustrating that substantial embedded compute is now
compatible with practical \gls{uav} \gls{swap} limits. These developments suggest that recent 
improvements in radar miniaturisation and embedded heterogeneous computing are promoting a 
new class of computation-heavy UAV sensing systems, making it timely to revisit radar 
target tracking and to determine the most suitable signal-processing pipeline for UAV 
detect-and-avoid applications. 

Despite the challenges, \gls{uav}s are the fastest growing market for radar systems, projected
to grow with a \gls{cagr} of 11.22\% globally through 2030 \cite{mordor_airborne_radars_market_2025}. To make these emerging radar systems more 
performant, accessible, and safe, it would be valuable to find the most suitable radar 
signal processing pipeline for UAV DAA systems, revisiting the old problem of radar target 
tracking in the face of unique hardware conditions.

Software-defined architecture (SDAs) are increasingly preferred over hardware defined
radar systems for \gls{uav} applications as software modifications can be shipped faster than
hardware modifications \cite{mordor_airborne_radars_market_2025}. As a result, \gls{gpu}s
are favoured over non-programmable digital signal processors such as \gls{asic}s.
\gls{gpu}-accelerated implementations are also favourable for computationally intensive
stages of the radar processing chain, as many core operations, such as \gls{fft}s and
matrix operations, are highly data-parallel, allowing for large increases in throughput
when compared to serial processors. Additionally, \gls{gpu}-based pipelines can leverage
mature developer friendly, \gls{gpu} software ecosystems and optimised libraries
\cite{zaknoun2024_realtime_gpu_accelerated_radar_processing}. Additionally,
the incorporation of AI into radar tracking algorithms has further encouraged these highly
parallel, \gls{gpu}-based pipelines \cite{electronics13214251}. Key open questions remain around
balancing the compute saved through \gls{gpu} parallelism against the cost of data movement and
synchronisation, and how unified (shared) \gls{cpu}--\gls{gpu} memory on embedded platforms changes
this trade-off.

Given the miniaturisation of radar enabling technology, the potential benefits for radar 
base DAA systems, and the range limitations of systems suitable for small radar systems. 
This report proposes a software pipeline design that allows for the radar tracking of 
aerial hazards on UAVs with strict SWaP constraints. Particular emphasis will be placed on 
maximising the sensitivity to faraway targets, a requirement for operating in shared 
airspace, and minimising computation resources and latency. 

Existing research demonstrates 
that compact radar hardware is now sufficiently mature for UAV DAA applications. 
SWaP-compliant systems have been built, detection at DAA-relevant ranges has been 
achieved, and processing feasibility on constrained embedded platforms has been 
established. However, three important gaps are identified.

First, existing hardware demonstrations validate detection only against large, 
high radar cross-section targets such as buildings and fixed-wing aircraft, 
leaving the practical sensitivity floor of compact UAV radar systems against 
small moving airborne targets unknown.

Second, no work has benchmarked a GPU-based radar pipeline in a UAV DAA 
setting. Existing studies are limited to FPGA and ARM CPU implementations, 
evaluate single algorithm sequences on single hardware configurations, and do 
not profile architectural bottlenecks such as host--device data transfer and 
thread synchronisation overhead.

Thirdly, angle sensing algorithms for low-SNR monopulse radar have only been 
validated under idealised white Gaussian noise with single simulated targets; 
no studies integrate these algorithms into a complete end-to-end range--Doppler 
tracking pipeline for UAV DAA. 

The core problem this thesis addresses is the absence of a principled, 
hardware-aware understanding of how algorithm selection and implementation 
choices jointly affect the sensitivity and computational performance of compact 
monopulse radar systems for UAV DAA. This thesis constructs that understanding 
through a systems-level investigation of the software design space of a 
GPU-based radar tracking pipeline, with specific contributions as follows:

\begin{itemize}
    \item A comparison of the algorithms necessary for monopulse radar, 
    assessed by detection sensitivity and track quality 
    in multi-target aerial scenarios, to inform algorithm selection.
    \item A comparison of alternative low-level implementations of the selected 
    algorithms, quantifying their effects on end-to-end throughput and latency 
    in highly parallelised GPU pipelines, to inform implementation choices.
    \item An assessment of how these algorithm and implementation conclusions 
    vary across embedded hardware platforms, providing hardware-aware 
    design recommendations.
\end{itemize}

This thesis concerns the software design and evaluation of radar tracking pipelines for UAV 
DAA under SWaP constraints. The scope covers processing stages from raw radar samples through
windowing, FFT, coherent integration, CFAR detection, monopulse angle estimation, data 
association, and tracking filters, on embedded GPU platforms suitable for small and medium UAVs.

Radar hardware design, full-scale field trials, FPGA and ASIC implementations, and fusion  
with optical or infrared sensing are outside scope. The thesis does not claim certification 
readiness; it provides comparative evidence and design guidance for onboard radar processing 
under realistic embedded constraints.

The radar sensor is assumed to be a small monopulse radar whose aperture and power budget are 
representative of a UAV-compatible installation. The target compute platform is GPU-enabled 
embedded hardware with shared or unified CPU/GPU memory, as found in many embedded systems-on-module, 
though CPU-only and discrete-memory configurations are included where needed to clarify architectural 
trade-offs.

By benchmarking alternative implementations of radar signal processing stages and comparing 
tracking algorithms, this work quantifies how these choices affect detection sensitivity, 
false alarm rate, and computational performance. In addition to characterising the algorithms 
themselves, the project produces a reproducible evaluation framework (including computational 
profiling, ablation studies, and parameter sweeps) that can be applied to comparable radar 
research. The resulting design recommendations provide an evidence-based guide for engineers 
developing small radar systems, and support the integration of UAVs into civil airspace where 
reliable onboard tracking of non-cooperative targets is essential.

More broadly, improved DAA capability supports the safe scaling of beyond-visual-line-of-sight 
UAV operations in shared airspace, enabling high-value public and commercial services. These 
include faster emergency response and situational awareness, routine environmental and 
agricultural monitoring, and cost-effective inspection of critical infrastructure such as 
power lines and transport corridors \cite{casa_bvlos_trial_2025}. Strengthening the safety 
case for UAV integration also supports national bushfire risk reduction and response, 
where drones are an identified tool for monitoring, assessment, and operations planning 
\cite{drones_gov_bushfire}. By reducing uncertainty in achievable performance under SWaP 
constraints, this work contributes evidence that can inform practical requirements and 
deployment decisions for safe UAV operations in shared airspace \cite{utm_action_plan_2024}.

\newpage
\section{Thesis Structure}
\textbf{Chapter 2 - Background} introduces the background information related 
to the following areas: standard radar algorithms, monopulse angle estimation, 
multi-target tracking, and parallel processing.

\textbf{Chapter 3 - Literature Review} conducts and extensive literature 
review of both foundational and state of the art approaches to performant 
radar systems under high SWaP constraints. The methodologies of works with 
similar motivations to ours are critically evaluated for their effectiveness 
in meeting their aims. Key findings and gaps across works are summarised and 
an explanation for how this thesis fits in with the existing literature is 
provided. 

\textbf{Chapter 4 - Methodology} Explains the methodology by which we compared 
different radar processing pipeline designs for their suitability for use in 
detect and avoid systems, and compare different algorithms for their ability 
to detect small, faraway targets. This methodology will describe the 
algorithms, the implementations of these algorithms, the processes by which 
the implementations are profiled, and the processes by which the algorithms 
are compared for their impact on the overall sensitivity of the system. 

\textbf{Chapter 5 - Results} Presents the results associated with each of the 
experiments outlined in the methodology. The results are illustrated as tables 
and figures, with descriptions of the findings of each of the methodology's 
stages.

\textbf{Chapter 6 - Discussion} The key findings in the results are discussed. 
The results of different experiment stages are evaluated and explained. 
Recommendations are provided. The results are tied in with the problem 
and contributions introduce in this chapter and Chapter 3. The contributions 
that were not achieved are reflected on. 

\textbf{Chapter 7 - Project Plan} A summary of what was achieved so far is provided, and 
a plan for the remaining work is outlined. The plan includes a timeline, stakeholder management plan, and 
a risk management plan.   
